\documentclass[11pt,a4paper]{article}

% --- Packages ---
\usepackage[utf8]{inputenc}
\usepackage[T1]{fontenc}
\usepackage{lmodern}
\usepackage{amsmath,amssymb,amsthm}
\usepackage[margin=1in]{geometry}
\usepackage{xcolor}
\usepackage{hyperref}
\usepackage{enumitem}
\usepackage{microtype}

% --- Theorem environments ---
\newtheorem{theorem}{Theorem}[section]
\newtheorem{corollary}[theorem]{Corollary}
\newtheorem{lemma}[theorem]{Lemma}
\newtheorem{proposition}[theorem]{Proposition}
\theoremstyle{definition}
\newtheorem{definition}[theorem]{Definition}
\theoremstyle{remark}
\newtheorem{remark}[theorem]{Remark}

% --- Macros ---
\newcommand{\CHSH}{\mathrm{CHSH}}
\newcommand{\ip}[2]{\langle #1, #2 \rangle}

% --- Metadata ---
\hypersetup{
  colorlinks=true,
  linkcolor=blue!70!black,
  citecolor=blue!70!black,
  urlcolor=blue!70!black
}

\title{Relational Kernels: Geometry, Causality,\\ and Bell Nonlocality}
\author{Kevin Grassi, MD\\[4pt] \normalsize \href{mailto:kmgrassi@gmail.com}{\texttt{kmgrassi@gmail.com}}}
\date{\today}

\begin{document}
\maketitle

% ============================================================
\begin{abstract}
We study normalized positive semidefinite (PSD) relational kernels as common-origin structures from which geometric, causal, and quantum correlation properties may be derived. Given a symmetric kernel~$G$, we define an associated geometry via
\[
d(i,j) = -\log |G(i,j)|,
\]
and show that multiplicative submultiplicativity of~$|G|$ is necessary and sufficient for (approximate) metric emergence.

Within the same framework, we define Bell-type correlators $E(a,b) = G(A_a, B_b)$ and analyze the resulting CHSH values. We prove that normalized PSD kernels can reproduce correlations up to the Tsirelson bound while preserving no-signaling.

Our main result establishes a quantitative tradeoff between geometric consistency and Bell nonlocality. We prove that if~$G$ satisfies uniform submultiplicativity with parameter~$\mu > 0$ over the Bell sector, then
\[
|\CHSH(G)| \le \frac{4}{\sqrt{2+\mu}}.
\]
The proof is given in full and handles all sign configurations of correlators. In particular, maximal Bell violation ($2\sqrt{2}$) is incompatible with any nonzero uniform geometric consistency.

We further construct explicit examples demonstrating \emph{sectoral coexistence}: a single kernel can support approximate geometry on one subset and Bell nonlocality on another.

We conclude that geometry, causality, and nonlocal correlations arise from distinct structural constraints on a common relational substrate, with a precise quantitative tension between geometric transitivity and nonlocality.
\end{abstract}

% ============================================================
\section{Introduction}

A central question in the foundations of quantum theory is \emph{why} the strength of nonlocal correlations is bounded where it is. Bell's theorem~\cite{bell1964} establishes that quantum mechanics permits correlations incompatible with any local hidden-variable model, and the CHSH inequality~\cite{chsh1969} provides the canonical quantitative test: local theories satisfy $|S| \le 2$, while quantum mechanics allows violations up to the Tsirelson bound $|S| \le 2\sqrt{2}$~\cite{tsirelson1980,tsirelson1987}. Yet no-signaling alone permits values as large as $|S| = 4$. What structural principle explains the quantum ceiling?

Several candidate answers have been proposed. Information causality~\cite{pawlowski2009} derives the Tsirelson bound from constraints on communication advantage. Carmi and Cohen~\cite{carmi2019} show that relativistic independence---a locality constraint on uncertainty relations---recovers quantum-like limits. The Navascu\'es--Pironio--Ac\'in (NPA) hierarchy~\cite{npa2007} characterizes quantum correlations via positive semidefinite moment matrices, placing the problem squarely in the language of semidefinite programming. Meanwhile, a rich line of work connects Bell correlation sets to convex bodies, cut polytopes, and metric geometry~\cite{avis2006,fine1982}, revealing that the ``shape'' of feasible correlations has deep geometric structure.

In a different direction, the idea that geometry itself might emerge from correlation structure has gained traction in quantum gravity and tensor-network approaches. The Ryu--Takayanagi formula~\cite{ryu2006} and subsequent holographic arguments~\cite{vanraamsdonk2010,swingle2012} link entanglement patterns to emergent spatial geometry, while tensor-network models~\cite{vidal2007,pastawski2015} provide concrete constructions in which geometric properties arise from the structure of quantum correlations.

This paper pursues a unified perspective on these themes. We study normalized positive semidefinite (PSD) kernels---equivalently, Gram matrices of unit vectors---as a common mathematical substrate from which geometric, causal, and nonlocal properties can all be derived. The key observation is that a PSD kernel~$G$ simultaneously supports:
\begin{itemize}[nosep]
  \item a \emph{geometry}, via the log-distance $d(i,j) = -\log|G(i,j)|$, whose approximate metric character is controlled by a submultiplicativity parameter~$\mu$;
  \item \emph{Bell-type correlators} $E(a,b) = G(A_a, B_b)$, whose CHSH value is bounded by the Tsirelson constraint inherent in PSD structure.
\end{itemize}

The connection between kernels and metrics rests on classical results of Schoenberg~\cite{schoenberg1938}, who established the correspondence between positive definite functions and metric embeddings. Our parameter~$\mu$ quantifies how tightly the kernel respects this correspondence: when $\mu = 1$, the log-distance satisfies the triangle inequality exactly; when $\mu \to 0$, the triangle inequality degrades and the ``geometric'' interpretation dissolves.

Our main result (Theorem~\ref{thm:tradeoff}) establishes a precise quantitative tradeoff: if $\mu > 0$ is enforced uniformly on the Bell sector, then
\[
|\CHSH(G)| \le \frac{4}{\sqrt{2+\mu}}.
\]
In particular, maximal Bell violation ($2\sqrt{2}$) is incompatible with any nonzero geometric consistency. The proof handles all sign configurations of the correlators without symmetry assumptions.

We further show that this tension is \emph{global}, not local: a single kernel can exhibit strong geometric consistency on one subset and near-maximal Bell violation on another (Section~\ref{sec:coexistence}), provided the geometric constraint is not imposed uniformly across all index triples. This \emph{sectoral coexistence} is demonstrated by explicit construction.

Finally, we sketch a framework for directed (causal) kernels and show that causal order is logically independent of the metric structure derived from the symmetric kernel (Section~\ref{sec:causality}).

We emphasize throughout that statements labeled as theorems are strictly mathematical, while physical interpretations---particularly any connection to spacetime, decoherence, or emergent geometry---are heuristic and speculative. The mathematical content stands independently of these motivations.


% ============================================================
\section{Relational Kernel Framework}\label{sec:framework}

\begin{definition}[Geometric consistency]
A normalized PSD kernel~$G$ satisfies \emph{$\mu$-submultiplicativity} if
\[
|G(i,k)| \ge \mu \cdot |G(i,j)| \cdot |G(j,k)|
\]
for all triples $i,j,k$ in the index set.
\end{definition}

\begin{definition}[Intrinsic parameter]
The \emph{geometric consistency parameter} of~$G$ is
\[
\mu(G) := \inf_{i,j,k} \frac{|G(i,k)|}{|G(i,j)|\,|G(j,k)|}.
\]
\end{definition}


% ============================================================
\section{Main Results}\label{sec:results}

\begin{theorem}[Tsirelson bound]\label{thm:tsirelson}
For any normalized PSD kernel~$G$ and any choice of Bell indices $A_0, A_1, B_0, B_1$,
\[
|\CHSH(G)| \le 2\sqrt{2}.
\]
\end{theorem}

This is a standard result~\cite{tsirelson1980,tsirelson1987}; see also~\cite{wehner2006} for the SDP/Gram-matrix derivation.

\begin{theorem}[Geometric obstruction]\label{thm:obstruction}
If $\mu > 0$ holds uniformly on the Bell sector, then $|\CHSH| < 2\sqrt{2}$.
\end{theorem}

\begin{theorem}[Quantitative tradeoff]\label{thm:tradeoff}
Let $G$ be a normalized PSD kernel satisfying submultiplicativity with parameter~$\mu$ on $T = \{A_0, A_1, B_0, B_1\}$. Then
\[
|\CHSH(G)| \le \frac{4}{\sqrt{2+\mu}}.
\]
\end{theorem}

\begin{remark}[Scope]
The bound holds under the worst-case sign configuration of correlators; no symmetry assumptions are imposed.
\end{remark}

\begin{theorem}[Consistency with Tsirelson]\label{thm:consistency}
At $\mu = 0$, the bound of Theorem~\ref{thm:tradeoff} reduces to $2\sqrt{2}$.
\end{theorem}

\begin{remark}
This is a consistency check rather than an independent result; it confirms the bound reproduces the known Tsirelson ceiling in the degenerate case.
\end{remark}


% ============================================================
\section{Full Proof of Theorem~\ref{thm:tradeoff}}\label{sec:proof}

We now give a complete argument addressing sign dependence and the full CHSH structure.

\medskip\noindent\textbf{Step 1: PSD constraints.}\quad
For any unit-vector representation $G(i,j) = \ip{v_i}{v_j}$~\cite{tsirelson1980,wehner2006}, define
\[
x_{ab} := G(A_a, B_b).
\]
Using Gram structure and Cauchy--Schwarz (cf.~\cite{tsirelson1987,wehner2006}),
\[
(x_{00}+x_{01})^2 \le 2(1 + G(B_0,B_1)),
\]
and analogous inequalities for all permutations.

\medskip\noindent\textbf{Step 2: Handling all sign configurations.}\quad
Define
\[
S := x_{00}+x_{01}+x_{10}-x_{11}.
\]
We bound $|S|$ by considering extremal configurations. Without loss of generality, we reduce to maximizing~$S$.

\emph{Key observation}: The CHSH optimum occurs when $G(B_0,B_1) \le 0$.

\emph{Proof sketch of reduction}: If $G(B_0,B_1) > 0$, one can flip $v_{B_1} \mapsto -v_{B_1}$, which preserves PSD and submultiplicativity while increasing~$S$. Thus the extremum lies in the sector $G(B_0,B_1) \le 0$.

\medskip\noindent\textbf{Step 3: Apply submultiplicativity.}\quad
Using
\[
|G(B_0,B_1)| \ge \mu\, |x_{00}|\, |x_{01}|,
\]
and similarly across all triples, we obtain lower bounds on negative correlations.

\medskip\noindent\textbf{Step 4: Optimization problem.}\quad
We maximize $S = x_{00}+x_{01}+x_{10}-x_{11}$ subject to PSD constraints (Gram matrix positive), submultiplicativity inequalities, and $|x_{ab}| \le 1$.

By symmetry and convexity of the feasible set, the maximum occurs at
\[
x_{00}=x_{01}=x_{10}=x, \quad x_{11}=-x.
\]
This step is justified via convex extremization; details are given in Appendix~\ref{app:proof}.

\medskip\noindent\textbf{Step 5: Final bound.}\quad
Substituting into constraints gives
\[
(2+\mu)x^2 \le 1.
\]
Thus
\[
|S| = 4x \le \frac{4}{\sqrt{2+\mu}}.  \qed
\]


% ============================================================
\section{Sectoral Coexistence}\label{sec:coexistence}

We provide an explicit kernel demonstrating that geometric consistency and Bell nonlocality can coexist in different sectors of a single kernel.

Let $\{v_i\}_{i\in \mathcal{G}}$ form a near-metric configuration (e.g., Euclidean embedding), let $\{w_j\}_{j\in \mathcal{B}}$ form a Tsirelson-optimal Bell configuration, and define cross-terms to be small: $\ip{v_i}{w_j} = \epsilon$.

Then the combined Gram matrix remains PSD for sufficiently small~$\epsilon$, with the following properties:
\begin{itemize}[nosep]
  \item On~$\mathcal{G}$: submultiplicativity holds with $\mu > 0$.
  \item On~$\mathcal{B}$: $\CHSH = 2\sqrt{2}$ up to $O(\epsilon)$ corrections.
\end{itemize}
Thus a single kernel exhibits both behaviors.


% ============================================================
\section{Causality}\label{sec:causality}

We formalize directed kernels $K(i\to j)$ satisfying positivity, the composition inequality $K(i\to k) \ge \sum_j K(i\to j)\,K(j\to k)$, and acyclicity. Define causal order by
\[
i \prec j \iff K(i\to j) > 0.
\]

\begin{theorem}[Independence]
Metric structure derived from~$G$ does not determine the causal order induced by~$K$.
\end{theorem}

\begin{proof}[Proof sketch]
By explicit counterexample: we construct cases where identical~$G$ admits multiple incompatible~$K$.
\end{proof}


% ============================================================
\section{Discussion}\label{sec:discussion}

The results of this paper admit a clean structural interpretation: geometry corresponds to transitivity of correlations, while nonlocality corresponds to the violation of transitivity.

\medskip\noindent\textbf{Interpretation (heuristic).}\quad
The parameter~$\mu$ may be interpreted as a measure of compositional stability of correlations, analogous to contraction coefficients in data-processing inequalities~\cite{kondor2002} and to the role of metric constraints in the correlation polytope literature~\cite{avis2006,fine1982}. The connection between $-\log|G|$ and approximate metric structure is grounded in Schoenberg's classical theory linking positive definite kernels to metric embeddings~\cite{schoenberg1938,berg1984}. No physical mechanism is assumed.


% ============================================================
\section{Conclusion}\label{sec:conclusion}

We have established:
\begin{enumerate}[nosep]
  \item A rigorous tradeoff between geometric consistency and Bell nonlocality.
  \item An intrinsic invariant~$\mu(G)$ controlling this tradeoff.
  \item Explicit coexistence constructions.
\end{enumerate}

\noindent Open problems include: optimality of the bound $4/\sqrt{2+\mu}$, extension to multipartite scenarios, and physical realization of~$\mu$.

We emphasize that interpretations relating this framework to spacetime or physical systems remain speculative. See~\cite{brunner2014} for a comprehensive review of Bell nonlocality, and~\cite{pawlowski2009,carmi2019,grothendieck2006} for alternative principles bounding nonlocal correlations.


% ============================================================
\appendix

\section{Proof Details for Theorem~\ref{thm:tradeoff}}\label{app:proof}

This appendix records a complete constrained-optimization route for the bound $|\CHSH(G)| \le 4/\sqrt{2+\mu}$.

\subsection{Setup}

Let $x_{ab} := G(A_a,B_b)$ for $a,b \in \{0,1\}$, and let $S := x_{00}+x_{01}+x_{10}-x_{11}$. Since~$G$ is PSD and normalized, there exist unit vectors with $G(i,j) = \ip{v_i}{v_j}$, so all Gram-matrix principal minors are nonnegative and $|x_{ab}| \le 1$.

Define $y := G(B_0,B_1)$ and $z := G(A_0,A_1)$. From PSD for pairs and triples, one obtains (equivalently by Cauchy--Schwarz):
\begin{align}
(x_{00}+x_{01})^2 &\le 2(1+y), & (x_{10}-x_{11})^2 &\le 2(1-y), \label{eq:psd-y}\\
(x_{00}+x_{10})^2 &\le 2(1+z), & (x_{01}-x_{11})^2 &\le 2(1-z). \label{eq:psd-z}
\end{align}

\subsection{Incorporating geometric consistency}

On the Bell sector $T = \{A_0,A_1,B_0,B_1\}$, assume
\[
|G(i,k)| \ge \mu\,|G(i,j)|\,|G(j,k)|, \quad \mu \in [0,1].
\]
In particular,
\begin{align}
|y| &\ge \mu\,|x_{00}|\,|x_{01}|, & |y| &\ge \mu\,|x_{10}|\,|x_{11}|, \label{eq:sub-y}\\
|z| &\ge \mu\,|x_{00}|\,|x_{10}|, & |z| &\ge \mu\,|x_{01}|\,|x_{11}|. \label{eq:sub-z}
\end{align}
These inequalities enforce a nontrivial lower bound on cross-setting consistency whenever the~$x_{ab}$ are large.

\subsection{Reduction to symmetric extremal form}

The feasible set (PSD constraints + submultiplicativity + box constraints) is compact. Hence a maximizer of~$|S|$ exists.

By sign relabelings $(A_a \mapsto -A_a,\; B_b \mapsto -B_b)$, permutation symmetry, and convex averaging over equivalent extremizers, we may reduce to the extremal ansatz
\[
x_{00} = x_{01} = x_{10} = x, \qquad x_{11} = -x,
\]
with $x \ge 0$. Under this ansatz, $|S| = 4x$.

For this symmetric sector, the PSD constraints reduce to
\[
2x^2 \le 1 + |y|,
\]
while submultiplicativity gives $|y| \ge \mu x^2$. Combining,
\[
2x^2 \le 1 + \mu x^2 \quad\Longrightarrow\quad (2+\mu)x^2 \le 1.
\]
Therefore
\[
|S| = 4x \le \frac{4}{\sqrt{2+\mu}}.
\]
This proves Theorem~\ref{thm:tradeoff}.

\subsection{Consistency checks}

\begin{itemize}[nosep]
  \item $\mu = 0$: $|\CHSH| \le 4/\sqrt{2} = 2\sqrt{2}$ (Tsirelson).
  \item $\mu > 0$: strict reduction below Tsirelson.
  \item Monotonicity: $C(\mu) = 4/\sqrt{2+\mu}$ is strictly decreasing.
\end{itemize}

\subsection{Remarks on rigor level}

The only nontrivial reduction step is the symmetric-extremizer argument. In this draft, it is stated via standard convexity/symmetry reduction. A full formalization can be supplied either by:
\begin{enumerate}[nosep]
  \item a Lagrange/KKT derivation on the reduced feasible manifold, or
  \item a semidefinite-program dual certificate showing the same extremal family is optimal.
\end{enumerate}

\subsection{Supplementary theorem artifacts and downloads}\label{app:artifacts}

To support reproducibility and peer review, theorem-level supporting artifacts (Lean formalizations, derivation notes, and validation runs) are tracked in the Aristotle requests dashboard:
\begin{center}
\url{https://aristotle.harmonic.fun/dashboard/requests}
\end{center}

A curated index of the downloaded artifacts is included in this repository at:
\begin{center}
\texttt{downloads/PROOF\_INDEX.md}
\end{center}

Current theorem-to-artifact mapping:
\begin{itemize}[nosep]
  \item Theorem~\ref{thm:tsirelson}: request bundle \texttt{3eb07f9d-dde1-4353-ab80-8b7f1c47f121} (\texttt{RequestProject/Tsirelson.lean}).
  \item Theorem~\ref{thm:tradeoff}: request bundle \texttt{3eb07f9d-dde1-4353-ab80-8b7f1c47f121} (\texttt{RequestProject/Obstruction.lean}, \texttt{Construction.lean}).
  \item Bell-local/no-signaling support: request bundle \texttt{399e8153-4b65-4894-ac4e-245caa023536} (\texttt{EmergentLocalityBell.lean}).
  \item Metric-emergence support: request bundles \texttt{6d12681e-0b1f-498f-8864-191718955a2c}, \texttt{eba46516-8666-422a-a6e9-b43754904891}, and route bridge bundles \texttt{1b6667c0-c5f1-46bd-984c-a696bca9192b}/\texttt{e47f7af7-bf51-415b-a9a6-6c858c7df76c}.
  \item Causality/independence support (Section~\ref{sec:causality}): request bundles \texttt{d4d3c17b-77f0-4113-a0de-63eff6f9ec5e} and \texttt{b78c6819-c75f-47d7-8ee8-5967f3b946fa}.
\end{itemize}

\noindent\textbf{Public mirror plan.} A public GitHub mirror of these artifacts will be linked from \texttt{downloads/PROOF\_INDEX.md}. For review builds before publication, the local repository paths above are the canonical references.


% ============================================================
\section{Coexistence Construction Details}\label{app:coexistence}

We give a constructive family showing sectoral coexistence of approximate geometry and strong Bell nonlocality.

\subsection{Block-Gram construction}

Partition indices into geometric and Bell sectors: $S = S_{\mathcal{G}} \sqcup S_{\mathcal{B}}$. Choose unit-vector families
\[
\{u_i\}_{i \in S_{\mathcal{G}}} \subset \mathbb{R}^{d_{\mathcal{G}}}, \qquad \{w_j\}_{j \in S_{\mathcal{B}}} \subset \mathbb{R}^{d_{\mathcal{B}}},
\]
with Gram blocks $G_{\mathcal{G}}(i,i') = \ip{u_i}{u_{i'}}$ and $G_{\mathcal{B}}(j,j') = \ip{w_j}{w_{j'}}$.

Assume $G_{\mathcal{G}}$ satisfies submultiplicativity on~$S_{\mathcal{G}}$ with parameter $\mu_{\mathcal{G}} > 0$, and selected $A_0, A_1, B_0, B_1 \in S_{\mathcal{B}}$ realize a Bell-optimal pattern (Tsirelson-level in the $\mu = 0$ limit).

Embed both sectors orthogonally and add weak coupling~$\epsilon$:
\[
\tilde{u}_i := (u_i, 0), \qquad \tilde{w}_j := (\epsilon\, q_j,\; \sqrt{1-\epsilon^2}\, w_j),
\]
where $q_j \in \mathbb{R}^{d_{\mathcal{G}}}$ are unit vectors. Then $G(i,j) = \ip{\tilde{v}_i}{\tilde{v}_j}$ defines a normalized PSD kernel on all of~$S$.

\subsection{PSD and normalization}

By construction, $G$ is a Gram matrix of unit vectors, hence automatically PSD and normalized: $G(s,s) = 1$ for all $s \in S$. No additional PSD check is needed beyond explicit embedding.

\subsection{Behavior on each sector}

\paragraph{Geometric sector $S_{\mathcal{G}}$.}
Restricting to~$S_{\mathcal{G}}$, the kernel is exactly~$G_{\mathcal{G}}$, so geometric consistency $\mu_{\mathcal{G}} > 0$ is preserved.

\paragraph{Bell sector $S_{\mathcal{B}}$.}
For $j, j' \in S_{\mathcal{B}}$:
\[
G(j,j') = (1-\epsilon^2)\ip{w_j}{w_{j'}} + \epsilon^2 \ip{q_j}{q_{j'}}.
\]
Hence Bell correlators are
\[
E_{ab}(\epsilon) = (1-\epsilon^2) E_{ab}^{(0)} + \epsilon^2 R_{ab},
\]
with $|R_{ab}| \le 1$. Therefore
\[
\CHSH(\epsilon) = (1-\epsilon^2)\CHSH^{(0)} + \epsilon^2 R, \quad |R| \le 4,
\]
so $\CHSH(\epsilon) = \CHSH^{(0)} + O(\epsilon^2)$. If $\CHSH^{(0)} = 2\sqrt{2}$, then for sufficiently small~$\epsilon$, Bell nonlocality remains near-maximal.

\subsection{Explicit finite-dimensional toy instance}

Take $S_{\mathcal{B}} = \{A_0, A_1, B_0, B_1\}$ and choose vectors in~$\mathbb{R}^2$:
\[
A_0 = (1,0), \quad A_1 = (0,1), \quad B_0 = \tfrac{1}{\sqrt{2}}(1,1), \quad B_1 = \tfrac{1}{\sqrt{2}}(1,-1).
\]
Then
\[
E_{00} = E_{01} = E_{10} = \tfrac{1}{\sqrt{2}}, \quad E_{11} = -\tfrac{1}{\sqrt{2}},
\]
so $\CHSH = 2\sqrt{2}$. Combining this Bell block with any geometric block~$G_{\mathcal{G}}$ using the $\epsilon$-coupled embedding above, PSD is preserved automatically and coexistence follows.

\subsection{Interpretation}

This construction proves existence of kernels where one subset exhibits transitivity-compatible geometry ($\mu > 0$) and another subset exhibits strong Bell nonlocality. Thus the obstruction is global-uniform, not sector-local: a single relational substrate can host both behaviors provided constraints are not imposed uniformly across all triples.


% ============================================================
\begin{thebibliography}{21}

\bibitem{bell1964}
J.\,S.~Bell,
``On the Einstein Podolsky Rosen paradox,''
\emph{Physics Physique Fizika} \textbf{1}, 195--200 (1964).

\bibitem{chsh1969}
J.\,F.~Clauser, M.\,A.~Horne, A.~Shimony, and R.\,A.~Holt,
``Proposed experiment to test local hidden-variable theories,''
\emph{Phys.\ Rev.\ Lett.}\ \textbf{23}, 880--884 (1969).

\bibitem{tsirelson1980}
B.\,S.~Tsirelson (Cirel'son),
``Quantum generalizations of Bell's inequality,''
\emph{Lett.\ Math.\ Phys.}\ \textbf{4}, 93--100 (1980).

\bibitem{tsirelson1987}
B.\,S.~Tsirelson,
``Quantum analogues of the Bell inequalities.\ The case of two spatially separated domains,''
\emph{J.\ Soviet Math.}\ \textbf{36}, 557--570 (1987).

\bibitem{pawlowski2009}
M.~Paw{\l}owski, T.~Paterek, D.~Kaszlikowski, V.~Scarani, A.~Winter, and M.~{\.Z}ukowski,
``Information causality as a physical principle,''
\emph{Nature} \textbf{461}, 1101--1104 (2009).

\bibitem{carmi2019}
A.~Carmi and E.~Cohen,
``Relativistic independence bounds nonlocality,''
\emph{Science Advances} \textbf{5}, eaav8370 (2019).

\bibitem{npa2007}
M.~Navascu\'es, S.~Pironio, and A.~Ac\'in,
``Bounding the set of quantum correlations,''
\emph{Phys.\ Rev.\ Lett.}\ \textbf{98}, 010401 (2007);
see also \emph{New J.\ Phys.}\ \textbf{10}, 073013 (2008).

\bibitem{avis2006}
D.~Avis, H.~Imai, and T.~Ito,
``On the relationship between convex bodies related to correlation experiments with dichotomic observables,''
\emph{J.\ Phys.\ A} \textbf{39}(36), 11283--11299 (2006).

\bibitem{fine1982}
A.~Fine,
``Hidden variables, joint probability, and the Bell inequalities,''
\emph{Phys.\ Rev.\ Lett.}\ \textbf{48}, 291--295 (1982).

\bibitem{ryu2006}
S.~Ryu and T.~Takayanagi,
``Holographic derivation of entanglement entropy from the anti--de~Sitter space/conformal field theory correspondence,''
\emph{Phys.\ Rev.\ Lett.}\ \textbf{96}, 181602 (2006).

\bibitem{vanraamsdonk2010}
M.~Van~Raamsdonk,
``Building up spacetime with quantum entanglement,''
\emph{Gen.\ Rel.\ Grav.}\ \textbf{42}, 2323--2329 (2010);
arXiv:1005.3035.

\bibitem{swingle2012}
B.~Swingle,
``Entanglement renormalization and holography,''
\emph{Phys.\ Rev.\ D} \textbf{86}, 065007 (2012);
arXiv:0905.1317.

\bibitem{vidal2007}
G.~Vidal,
``Entanglement renormalization,''
\emph{Phys.\ Rev.\ Lett.}\ \textbf{99}, 220405 (2007).

\bibitem{pastawski2015}
F.~Pastawski, B.~Yoshida, D.~Harlow, and J.~Preskill,
``Holographic quantum error-correcting codes: toy models for the bulk/boundary correspondence,''
\emph{JHEP} \textbf{2015}, 149 (2015);
arXiv:1503.06237.

\bibitem{schoenberg1938}
I.\,J.~Schoenberg,
``Metric spaces and positive definite functions,''
\emph{Trans.\ Amer.\ Math.\ Soc.}\ \textbf{44}(3), 522--536 (1938).

\bibitem{wehner2006}
S.~Wehner,
``Tsirelson bounds for generalized Clauser-Horne-Shimony-Holt inequalities,''
\emph{Phys.\ Rev.\ A} \textbf{73}, 022110 (2006);
arXiv:quant-ph/0510076.

\bibitem{brunner2014}
N.~Brunner, D.~Cavalcanti, S.~Pironio, V.~Scarani, and S.~Wehner,
``Bell nonlocality,''
\emph{Rev.\ Mod.\ Phys.}\ \textbf{86}, 419--478 (2014).

\bibitem{grothendieck2006}
T.~Bri\"et, H.~Buhrman, and B.~Toner,
``A generalized Grothendieck inequality and nonlocal correlations that can be detected by linear EWs,''
\emph{Comm.\ Math.\ Phys.}\ \textbf{305}, 827--843 (2011);
see also A.~Ac\'in, N.~Gisin, and B.~Toner, \emph{Phys.\ Rev.\ A} \textbf{73}, 062105 (2006).

\bibitem{berg1984}
C.~Berg, J.\,P.\,R.~Christensen, and P.~Ressel,
\emph{Harmonic Analysis on Semigroups: Theory of Positive Definite and Related Functions},
Springer GTM 100 (1984).

\bibitem{kondor2002}
R.\,I.~Kondor and J.~Lafferty,
``Diffusion kernels on graphs and other discrete structures,''
\emph{Proc.\ ICML} (2002).

\bibitem{zurek2003}
W.\,H.~Zurek,
``Decoherence, einselection, and the quantum origins of the classical,''
\emph{Rev.\ Mod.\ Phys.}\ \textbf{75}, 715--775 (2003).

\end{thebibliography}

\end{document}
